% !TeX spellcheck = es_ES
%\documentclass[letterpaper,12pt]{article}
\documentclass[letterpaper,11pt]{article}

% Language & Input
  \usepackage[english]{babel} % Language: English
  \usepackage[ansinew]{inputenc} % Input
% Font
  \usepackage[T1]{fontenc} % Font encoding
  \usepackage{lmodern,microtype} % Typeface
% Math
  \usepackage{amsmath,amsthm,amssymb,amsfonts} % AMS math
  \usepackage{dsfont,mathrsfs,ushort} % Math style
% Style
  \usepackage[bf,sf,tiny,center]{titlesec} % Section titles 
  \usepackage{geometry} % Page & margins
  \usepackage{setspace} % Spacing
  \usepackage{enumitem} % Lists
  \usepackage{abstract} % Abstract
  \usepackage{color}
% References
  \usepackage[longnamesfirst]{natbib} % Manages bibiography
% Figures
  \usepackage{graphicx} % Additional options for includefigure
  \usepackage[font={small}]{caption}
% Tables
  \usepackage{booktabs} % Table format
  \usepackage{dcolumn} % Align on the decimal point of numbers in tabular columns.
  \usepackage{multirow} % Multirows in a table
  \usepackage{lscape} % Tables in landscape form
  \usepackage{ctable} % Tables with captions and notes all the same width 
% Hyperlinks
  \usepackage{hyperref}
% Addons
  \usepackage[l2tabu, orthodox]{nag} % Warn the user about the use of such obso­lete things

%%%%%%%%%%%%%%%%%%%%%%%%%%%%%%% Style & Structure
%% Titles (Carlos Perez)
%  \titlelabel{\thetitle.}
%  \titleformat*{\subsection}{\centering\sffamily\itshape}
%  \titleformat{\subsubsection}[runin]{\sffamily}{}{0em}{}[.]
%  \titlespacing{\subsubsection}{0pt}{*2}{\wordsep}
  
% Title page (Luis Cabral)
\makeatletter
    \renewcommand{\@maketitle}{%
    \newpage\parindent0in\null\vskip2em%
    {\LARGE\textbf{\textsf{\@title}}\par\vskip2em}%
    \@author\par\vskip2em%
    \@date%
    \par
    \thispagestyle{empty}\setcounter{page}{1}%\newpage
    }%
\makeatother
\renewenvironment{abstract}% Change the way abstract is presented
{\vskip2em\noindent{\textbf{\textsf{Abstract.}}}\hspace*{1ex}}{}

% Sections
  \renewcommand\thesection{\arabic{section}}
  \renewcommand\thesubsection{\Alph{subsection}}
% Bibliography
  \bibliographystyle{apa} %Alternative: {plain} {aea} {econometrica}
% Spacing & lists and margins
  \geometry{margin=2.54cm}
  \setstretch{1.15} % Set spacing
% Color
  \definecolor{ntr}{rgb}{0.8,0,0} % not too red
% Abstract
  \renewcommand{\abstractnamefont}{\sffamily\small\bfseries}
  \setlength{\abstitleskip}{-1em}
% Hyperlinks
  \hypersetup{colorlinks=true,linkcolor=ntr,citecolor=ntr,urlcolor=black,bookmarksnumbered=true}
% Notes to tables and figures
% Significant stars in tables
  %\newcommand{\sym}[1]{\rlap{#1}}
  \def\sym#1{\ifmmode^{#1}\else\(^{#1}\)\fi}
% Games
% Miscellaneous

%%%%%%%%%%%%%%%%%%%%%%%%%%%%%%% Theorems
  \theoremstyle{plain}
  \newtheorem{prop}{Proposition}
  \newtheorem{assumption} {Assumption}
  \newtheorem{theorem}{Theorem}
  \newtheorem{lem}{Lemma}
  \newtheorem{corollary}{Corollary}
  \theoremstyle{definition}
  \newtheorem{defi}{Definition}
  \newtheorem{ex}{Example}
  \theoremstyle{remark}
  \newtheorem{remark}{Remark}
  \def\propautorefname{Proposition}
  \def\propositionautorefname{Proposition}
  \def\defiautorefname{Definition}
  \def\lemautorefname{Lemma}
  \def\exampleautorefname{Example}
  \def\corollaryautorefname{Corollary}
  \def\exautorefname{Example}

\newcommand{\SR}{\mathbb{R}}

\DeclareMathOperator*{\Max}{m{a}x}
\DeclareMathOperator*{\Min}{m{i}n}
\DeclareMathOperator{\rank}{rank}
\DeclareMathOperator{\plim}{plim}
\DeclareMathOperator{\Var}{Var}
\DeclareMathOperator{\Avar}{AVar}
\DeclareMathOperator{\Cov}{Cov}
\DeclareMathOperator{\E}{\mathbb{E}}
\newcommand{\pto}{\stackrel{p}{\longrightarrow}}
\newcommand{\dto}{\stackrel{d}{\longrightarrow}}
\newcommand{\adistr}{\stackrel{a}{\sim}}
\newcommand{\epsi}{\varepsilon}
\newcommand{\tmle}{\hat \theta_{ML}}
\newcommand{\tgmm}{\hat \theta_{GMM}}
\newcommand\independent{\protect\mathpalette{\protect\independenT}{\perp}}
\def\independenT#1#2{\mathrel{\rlap{$#1#2$}\mkern2mu{#1#2}}}
\DeclareMathOperator{\N}{Normal}

\DeclareMathOperator{\Corr}{Corr}

\begin{document}

	\title {An\'alisis Econom\'etrico  \\ Tarea 3}
	\author{Mag\'ister en Econom\'ia, Universidad Alberto Hurtado}
	\date{Fecha de entrega: 22 de octubre de 2026 (al principio de la ayudant\'ia)}

\maketitle

Total: 100 puntos

Los ejercicios 3, 4 y 5 requieren estimar modelos con datos. Utilice R (u otro software
estad\'istico) y adjunte su c\'odigo. Reporte errores est\'andar robustos a
heteroscedasticidad. Los datos est\'an en el repositorio del curso, en formato R
(\href{https://github.com/ramirode/econometria_mae/raw/main/problem-sets/datos/wage2.rds}{\texttt{wage2.rds}},
\href{https://github.com/ramirode/econometria_mae/raw/main/problem-sets/datos/hprice2.rds}{\texttt{hprice2.rds}}; se leen con \texttt{readRDS})
y en formato Stata
(\href{https://github.com/ramirode/econometria_mae/raw/main/problem-sets/datos/wage2.dta}{\texttt{wage2.dta}},
\href{https://github.com/ramirode/econometria_mae/raw/main/problem-sets/datos/hprice2.dta}{\texttt{hprice2.dta}}).

\begin{enumerate}

    \item (20 puntos) Suponga homoscedasticidad. En clases vimos que, en ese caso, la varianza
        asint\'otica del estimador MCO es $\Avar(\hat \beta)=\frac{\sigma^2}{N}\E(x_ix_i')^{-1}$.
        \begin{enumerate}
            \item (10 puntos) Considere el modelo de regresi\'on simple
                \begin{align*}
                y_i=\beta_0+\beta_1 x_{1i}+u_i,
                \end{align*}
                donde $\E(u_i)=0$, $\E(x_{1i}\, u_i)=0$ y $\E(u_i^2|x_{1i})=\sigma^2$. Utilice la
                f\'ormula anterior con $x_i=(1, x_{1i})'$ para demostrar que
                \begin{align*}
                \Avar(\hat \beta_1)=\frac{\sigma^2}{N}\frac{1}{\Var(x_{1i})}.
                \end{align*}
            \item (5 puntos) Considere ahora el modelo con dos variables explicativas
                \begin{align*}
                y_i=\beta_0+\beta_1 x_{1i}+\beta_2 x_{2i}+v_i,
                \end{align*}
                donde $\E(v_i)=\E(x_{1i}\,v_i)=\E(x_{2i}\,v_i)=0$ y
                $\E(v_i^2|x_{1i},x_{2i})=\sigma_v^2$. En clases vimos que
                \begin{align*}
                \Avar(\hat \beta_1)=\frac{\sigma_v^2}{N}\frac{1}{\Var(r_{1i})},
                \end{align*}
                donde $\gamma_0+\gamma_1x_{2i}$ es el proyector lineal de $x_{1i}$ en una
                constante y $x_{2i}$, y $r_{1i}=x_{1i}-\gamma_0-\gamma_1x_{2i}$ es su residuo.
                Demuestre que
                \begin{align*}
                \Var(r_{1i})=\Var(x_{1i})\,(1-\rho^2),
                \end{align*}
                donde $\rho$ es el coeficiente de correlaci\'on entre $x_{1i}$ y $x_{2i}$: la
                covarianza dividida por el producto de las desviaciones est\'andar,
                \begin{align*}
                \rho=\Corr(x_{1i},x_{2i})=\frac{\Cov(x_{1i},x_{2i})}{\sqrt{\Var(x_{1i})}\,\sqrt{\Var(x_{2i})}}.
                \end{align*}
            \item (5 puntos) Explique c\'omo depende $\Avar(\hat \beta_1)$ del tama\~no de la
                muestra, de la varianza del error, de la varianza de $x_1$ y de la correlaci\'on
                entre $x_1$ y $x_2$. ?`Qu\'e ocurre cuando $|\rho|\to 1$?
        \end{enumerate}

    \item (10 puntos) Un estudio reciente ha encontrado que la tasa de mortalidad de las
        personas que duermen 7 horas por noche es menor que la tasa de mortalidad de
        aquellas personas que duermen 8 o m\'as horas. El estudio utiliz\'o 1,1 millones de
        observaciones de una muestra aleatoria de chilenos entre 30 y 102 a\~nos de edad. Cada
        persona en la muestra fue seguida por 4 a\~nos. La tasa de mortalidad para personas que
        duermen 7 horas fue calculada como el n\'umero de muertes durante los 4 a\~nos entre las
        personas que duermen 7 horas sobre el n\'umero de personas que duermen 7 horas en la
        muestra. C\'alculos similares se realizaron para el resto de las categor\'ias. En base a
        los resultados del estudio, ?`usted recomendar\'ia a los chilenos que duermen 9 horas por
        noche reducir las horas de sue\~no a 7 horas por noche si quieren prolongar sus
        vidas? ?`Por qu\'e o por qu\'e no? Explique.

    \item (20 puntos) Utilice los datos de \texttt{wage2} del trabajo de Blackburn y
        Neumark (QJE, 1992): 935 hombres de entre 28 y 38 a\~nos en 1980. Las variables
        \texttt{wage} (salario mensual, en d\'olares), \texttt{married} (igual a 1 si est\'a
        casado) y \texttt{urban} (igual a 1 si vive en una zona urbana) definen el modelo
        \begin{align*}
        \log(wage_i)=\beta_0+\beta_1 \, married_i+\beta_2 \, urban_i+\beta_3 \, married_i\times urban_i+u_i.
        \end{align*}
        \begin{enumerate}
            \item (10 puntos) Estime el modelo. Escriba $\E(\log(wage)|married,urban)$ para
                los cuatro grupos (casado urbano, casado rural, soltero urbano y soltero rural)
                en funci\'on de los $\beta$'s, e interprete cada uno de los coeficientes
                estimados.
            \item (5 puntos) Calcule la media muestral de $\log(wage)$ en cada grupo y
                verifique que coincide con el valor predicho por el modelo. Explique por qu\'e.
            \item (5 puntos) Suponga que construye una dummy para cada uno de los cuatro grupos
                y estima la regresi\'on de $\log(wage)$ en las cuatro dummies y una constante.
                ?`Qu\'e ocurre? Explique.
        \end{enumerate}

    \item (30 puntos) Utilice los datos de \texttt{wage2} para estudiar los retornos de
        la educaci\'on. Adem\'as de las variables del ejercicio anterior, la base incluye
        \texttt{educ} (a\~nos de educaci\'on), \texttt{IQ} (resultado en un test de
        inteligencia), \texttt{exper} (a\~nos de experiencia laboral), \texttt{tenure} (a\~nos
        en el empleo actual) y \texttt{south} (igual a 1 si vive en el sur de EE.UU.).
        \begin{enumerate}
            \item (5 puntos) Utilice un modelo de regresi\'on simple log-nivel para estimar el
                efecto de un a\~no m\'as de educaci\'on ($educ$) en salarios ($wage$). Interprete
                los resultados. ?`Usted piensa que el coeficiente estimado mide una relaci\'on
                causal? Explique.
            \item (10 puntos) El modelo de regresi\'on simple omite, entre otras cosas, la
                habilidad innata. En los datos existe una proxy de la habilidad innata: $IQ$.
                Antes de estimar, prediga si el coeficiente de $educ$ aumentar\'a o
                disminuir\'a al incluir $IQ$ como variable explicativa adicional, y explique por
                qu\'e. Luego estime el modelo con $IQ$, interprete el coeficiente de $educ$ y
                compare con (a).
            \item (10 puntos) La f\'ormula del sesgo de variable omitida tambi\'en se cumple en
                la muestra, reemplazando momentos poblacionales por momentos muestrales. Primero,
                compruebe este resultado en los datos: regrese $IQ$ en $educ$ y, en base a esa
                estimaci\'on y a lo obtenido en (b), recupere la pendiente estimada en (a).
                Segundo, demuestre anal\'iticamente que este resultado se cumple siempre en la
                muestra.
            \item (5 puntos) Estime el efecto de la educaci\'on sobre los salarios controlando
                por $IQ$, $exper$, $tenure$, $married$, $south$ y $urban$. Compare con los
                resultados de (a) y (b).
        \end{enumerate}

    \item (20 puntos) Utilice los datos de \texttt{hprice2} del trabajo de Harrison y
        Rubinfeld (JEEM, 1978) sobre el efecto de la contaminaci\'on en el precio de las
        viviendas: 506 comunidades del \'area metropolitana de Boston en 1970. Las variables son
        \texttt{price} (precio mediano de las viviendas, en d\'olares), \texttt{nox}
        (concentraci\'on de \'oxido de nitr\'ogeno en el aire, en partes por 100 millones),
        \texttt{rooms} (n\'umero promedio de habitaciones), \texttt{dist} (distancia ponderada
        a cinco centros de empleo) y \texttt{stratio} (ratio estudiantes por
        profesor).
        \begin{enumerate}
            \item (5 puntos) Utilice un modelo de regresi\'on simple log-log para estimar el
                efecto de un aumento en $nox$ en $price$. Interprete los resultados. ?`Usted
                piensa que el coeficiente estimado mide una relaci\'on causal? Explique.
            \item (5 puntos) Agregue las siguientes variables explicativas al modelo:
                $rooms$, $\log(dist)$ y $stratio$. Interprete la estimaci\'on del efecto de un
                aumento de la contaminaci\'on en el precio de las viviendas, compare con (a) y
                explique la diferencia.
            \item (5 puntos) Obtenga la estimaci\'on del coeficiente de $\log(nox)$ del modelo
                en (b) utilizando la regresi\'on particionada. Primero, guarde los residuos de una
                regresi\'on de $\log(nox)$ en $rooms$, $\log(dist)$ y $stratio$. Segundo, regrese
                $\log(price)$ en esos residuos. Verifique que la estimaci\'on coincide con (b).
                No se preocupe por los errores est\'andar, s\'olo por las estimaciones puntuales.
            \item (5 puntos) Compare el error est\'andar del coeficiente de $\log(nox)$ en (a) y
                en (b). Calcule la correlaci\'on muestral entre $\log(nox)$ y $\log(dist)$ y
                utilice la intuici\'on del ejercicio 1 para explicar la diferencia.
        \end{enumerate}

\end{enumerate}
\end{document}
